% GENERATED FILE. Edit canonical lessons, then run npm run generate:books.
% canonical-source-hash: fnv1a64:9348e9ed12a32c15
% canonical-lessons: JA-C374-senro, JA-C374-fumikiri, JA-C374-teisha, JA-C374-hassha, JA-C374-chikayoru

\chapter{Along the Railway Line}
\label{ch:ja-along-the-railway-line}
\hlchaptermodality{374}

\begin{chapteropening}
\textbf{By the end of this chapter:} \emph{I can talk about trains stopping and leaving.}

Complete the last lesson of chapter 374: I can talk about trains stopping and leaving.
\end{chapteropening}

\section[\texorpdfstring{\ja{せんろ} (senro)}{せんろ (senro)}]{\ja{せんろ} (senro) — a railway track}

\label{lesson:JA-C374-senro}



\begin{quote}
Before the new one: say the Japanese for fare adjustment, then the Japanese for to pass by.
\end{quote}

\subsection*{You'll want to know: \ja{せんろ}}
\textbf{\ja{せんろ}} — \emph{senro} — \textquotedblleft{}a railway track\textquotedblright{}.

\begin{quote}
\textbf{\ja{せんろの} \ja{むこうに} \ja{こうえんが} \ja{あります。}} — \emph{senro no mukō ni kōen ga arimasu} — \textquotedblleft{}There is a park on the other side of the railway line.\textquotedblright{}
\end{quote}

\begin{cousinweb}
\textbf{\ja{せんろ}} (\emph{senro}) is Sino-Japanese (kango): two Chinese characters read with their on'yomi, readings borrowed from Chinese. They are \emph{sen}, \textquotedblleft{}line\textquotedblright{} (Middle Chinese \emph{sjenH}) and \emph{ro}, \textquotedblleft{}road\textquotedblright{} (\emph{luH}) (KANJIDIC2; Middle Chinese in Baxter's transcription).
\end{cousinweb}

\begin{grammarlens}[title={What you've built}]
One more word for getting around.
\end{grammarlens}

\subsection*{Your turn}
\begin{itemize}
\raggedright
  \item \emph{Say it:} \emph{senro}
  \item \emph{Say it:} \emph{senro}, once more
  \item \emph{Say it:} \emph{senro no mukō ni kōen ga arimasu} — There is a park on the other side of the railway line.
  \item \emph{Recall:} say the Japanese for fare adjustment, then the Japanese for to pass by, then say \emph{senro} again
\end{itemize}

\begin{tcolorbox}[breakable,colback=teal!4,colframe=teal!35!black,title={Before you move on}]
What does \ja{せんろ} mean? (\textquotedblleft{}A railway track\textquotedblright{}.) Say it once more.
\end{tcolorbox}
\section[\texorpdfstring{\ja{ふみきり} (fumikiri)}{ふみきり (fumikiri)}]{\ja{ふみきり} (fumikiri) — a level crossing}

\label{lesson:JA-C374-fumikiri}



\begin{quote}
Before the new one: say the Japanese for to pass by, then the Japanese for a railway track.
\end{quote}

\subsection*{You'll want to know: \ja{ふみきり}}
\textbf{\ja{ふみきり}} — \emph{fumikiri} — \textquotedblleft{}a level crossing\textquotedblright{}.

\begin{quote}
\textbf{\ja{ふみきりの} \ja{まえで} \ja{くるまが} \ja{とまりました。}} — \emph{fumikiri no mae de kuruma ga tomarimashita} — \textquotedblleft{}The car stopped in front of the level crossing.\textquotedblright{}
\end{quote}

\begin{cousinweb}
\textbf{\ja{ふみきり}} (\emph{fumikiri}) is a native Japanese word (wago), inherited rather than borrowed from Chinese, built from \emph{fumi}, \textquotedblleft{}stepping\textquotedblright{} (from \emph{fumu}) and \emph{kiri}, \textquotedblleft{}cutting\textquotedblright{} (from \emph{kiru}); its kanji were attached later as kun'yomi (JMdict).
\end{cousinweb}

\begin{grammarlens}[title={What you've built}]
One more word for getting around.
\end{grammarlens}

\subsection*{Your turn}
\begin{itemize}
\raggedright
  \item \emph{Say it:} \emph{fumikiri}
  \item \emph{Say it:} \emph{fumikiri}, once more
  \item \emph{Say it:} \emph{fumikiri no mae de kuruma ga tomarimashita} — The car stopped in front of the level crossing.
  \item \emph{Recall:} say the Japanese for to pass by, then the Japanese for a railway track, then say \emph{fumikiri} again
\end{itemize}

\begin{tcolorbox}[breakable,colback=teal!4,colframe=teal!35!black,title={Before you move on}]
What does \ja{ふみきり} mean? (\textquotedblleft{}A level crossing\textquotedblright{}.) Say it once more.
\end{tcolorbox}
\section[\texorpdfstring{\ja{ていしゃ} (teisha)}{ていしゃ (teisha)}]{\ja{ていしゃ} (teisha) — stopping (of a train or bus)}

\label{lesson:JA-C374-teisha}



\begin{quote}
Before the new one: say the Japanese for a railway track, then the Japanese for a level crossing.
\end{quote}

\subsection*{You'll want to know: \ja{ていしゃ}}
\textbf{\ja{ていしゃ}} — \emph{teisha} — \textquotedblleft{}stopping (of a train or bus)\textquotedblright{}.

\begin{quote}
\textbf{\ja{この} \ja{でんしゃは} \ja{つぎの} \ja{えきで} \ja{ていしゃします。}} — \emph{kono densha wa tsugi no eki de teisha shimasu} — \textquotedblleft{}This train stops at the next station.\textquotedblright{}
\end{quote}

\begin{cousinweb}
\textbf{\ja{ていしゃ}} (\emph{teisha}) is Sino-Japanese (kango): two Chinese characters read with their on'yomi, readings borrowed from Chinese. They are \emph{tei}, \textquotedblleft{}halt\textquotedblright{} (Middle Chinese \emph{deng}) and \emph{sha}, \textquotedblleft{}car\textquotedblright{} (\emph{tsyhae}) (KANJIDIC2; Middle Chinese in Baxter's transcription).
\end{cousinweb}

\begin{grammarlens}[title={What you've built}]
One more word for getting around.
\end{grammarlens}

\subsection*{Your turn}
\begin{itemize}
\raggedright
  \item \emph{Say it:} \emph{teisha}
  \item \emph{Say it:} \emph{teisha}, once more
  \item \emph{Say it:} \emph{kono densha wa tsugi no eki de teisha shimasu} — This train stops at the next station.
  \item \emph{Recall:} say the Japanese for a railway track, then the Japanese for a level crossing, then say \emph{teisha} again
\end{itemize}

\begin{tcolorbox}[breakable,colback=teal!4,colframe=teal!35!black,title={Before you move on}]
What does \ja{ていしゃ} mean? (\textquotedblleft{}Stopping (of a train or bus)\textquotedblright{}.) Say it once more.
\end{tcolorbox}
\section[\texorpdfstring{\ja{はっしゃ} (hassha)}{はっしゃ (hassha)}]{\ja{はっしゃ} (hassha) — departure (of a train or bus)}

\label{lesson:JA-C374-hassha}



\begin{quote}
Before the new one: say the Japanese for a level crossing, then the Japanese for stopping (of a train or bus).
\end{quote}

\subsection*{You'll want to know: \ja{はっしゃ}}
\textbf{\ja{はっしゃ}} — \emph{hassha} — \textquotedblleft{}departure (of a train or bus)\textquotedblright{}.

\begin{quote}
\textbf{\ja{でんしゃの} \ja{はっしゃは} \ja{くじです。}} — \emph{densha no hassha wa kuji desu} — \textquotedblleft{}The train leaves at nine.\textquotedblright{}
\end{quote}

\begin{cousinweb}
\textbf{\ja{はっしゃ}} (\emph{hassha}) is Sino-Japanese (kango): two Chinese characters read with their on'yomi, readings borrowed from Chinese. They are \emph{hatsu}, \textquotedblleft{}departure\textquotedblright{} (Middle Chinese \emph{pjot}) and \emph{sha}, \textquotedblleft{}car\textquotedblright{} (\emph{tsyhae}), joined as \emph{hassha} (KANJIDIC2; Middle Chinese in Baxter's transcription).
\end{cousinweb}

\begin{grammarlens}[title={What you've built}]
One more word for getting around.
\end{grammarlens}

\subsection*{Your turn}
\begin{itemize}
\raggedright
  \item \emph{Say it:} \emph{hassha}
  \item \emph{Say it:} \emph{hassha}, once more
  \item \emph{Say it:} \emph{densha no hassha wa kuji desu} — The train leaves at nine.
  \item \emph{Recall:} say the Japanese for a level crossing, then the Japanese for stopping (of a train or bus), then say \emph{hassha} again
\end{itemize}

\begin{tcolorbox}[breakable,colback=teal!4,colframe=teal!35!black,title={Before you move on}]
What does \ja{はっしゃ} mean? (\textquotedblleft{}Departure (of a train or bus)\textquotedblright{}.) Say it once more.
\end{tcolorbox}
\section[\texorpdfstring{\ja{ちかよる} (chikayoru)}{ちかよる (chikayoru)}]{\ja{ちかよる} (chikayoru) — to go near; to approach}

\label{lesson:JA-C374-chikayoru}



\begin{quote}
Before the new one: say the Japanese for stopping (of a train or bus), then the Japanese for departure (of a train or bus).
\end{quote}

\subsection*{You'll want to know: \ja{ちかよる}}
\textbf{\ja{ちかよる}} — \emph{chikayoru} — \textquotedblleft{}to go near; to approach\textquotedblright{}.

\begin{quote}
\textbf{\ja{せんろに} \ja{ちかよらないで} \ja{ください。}} — \emph{senro ni chikayoranaide kudasai} — \textquotedblleft{}Please keep away from the railway line.\textquotedblright{}
\end{quote}

Polite \textbf{\ja{ちかよります}}, \emph{chikayorimasu}; a godan verb. Past (polite): \textbf{\ja{こどもが} \ja{せんろに} \ja{ちかよりました。}}, \emph{kodomo ga senro ni chikayorimashita}.

\begin{cousinweb}
\textbf{\ja{ちかよる}} (\emph{chikayoru}) is a native Japanese word (wago), inherited rather than borrowed from Chinese, built from \emph{chika}, \textquotedblleft{}near\textquotedblright{} (from \emph{chikai}) and \emph{yoru}, \textquotedblleft{}to draw close\textquotedblright{}; its kanji were attached later as kun'yomi (JMdict).
\end{cousinweb}

\begin{grammarlens}[title={What you've built}]
One more word for getting around.
\end{grammarlens}

\subsection*{Your turn}
\begin{itemize}
\raggedright
  \item \emph{Say it:} \emph{chikayoru}
  \item \emph{Say it:} \emph{chikayoru}, once more
  \item \emph{Say it:} \emph{senro ni chikayoranaide kudasai} — Please keep away from the railway line.
  \item \emph{Recall:} say the Japanese for stopping (of a train or bus), then the Japanese for departure (of a train or bus), then say \emph{chikayoru} again
\end{itemize}

\begin{tcolorbox}[breakable,colback=teal!4,colframe=teal!35!black,title={Before you move on}]
What does \ja{ちかよる} mean? (\textquotedblleft{}To go near; to approach\textquotedblright{}.) Say it once more.
\end{tcolorbox}
